% Separate post-CP18 candidate. No unrestricted least-index conclusion.
\subsection{Residual irreducibility and a cycle of tame character types}

In this subsection ``irreducible residual inertia'' means irreducible
over $\mathbf F_p$ after passing to the completed maximal unramified
base. It does not mean absolutely irreducible over
$\overline{\mathbf F}_p$. The distinction matters: the irreducible
two-dimensional residual representations occurring here are tame
and split over $\mathbf F_{p^2}$.

\begin{lemma}[The leading-coordinate character]
\label{lem:B-RD53-leading}
\status{Proved; full independent proof and exact-code audits passed.}
Let $p\ge5$, and let $K$ be complete discretely valued of mixed
characteristic $(0,p)$ with algebraically closed residue field,
$e_K=p-1$, and $\mu_p\subset K$. Put $N=p^2-1$.
Fix a uniformizer of $K$, the associated compatible tame
fundamental characters, and a generator $\tau$ of tame inertia
whose level-two value is a primitive root $\zeta\in\mu_N$.
Let $M$ be potentially supersingular of height two, and assume
that $M[p]$ is irreducible over $\mathbf F_p$.
If the special tangent character takes the value
$\zeta^c$ at $\tau$, then the residual eigenvalues are
\[
              \bigl\{\zeta^{c+p-1},
                       \zeta^{p(c+p-1)}\bigr\}.       \tag{B-RD53.1}
\]
Here tangent means the action on the one-dimensional Lie space
of the special formal group. The special covariant Dieudonne
character has eigenvalues $\{\zeta^c,\zeta^{pc}\}$.
\end{lemma}
\begin{proof}
The image of wild inertia in $\GL(M[p])$ is a normal finite
$p$-group. Such a group has a nonzero invariant vector: its
orbits on the finite set $M[p]$ show that the number of fixed
vectors is divisible by $p$, and zero is fixed. Its invariant
subspace is stable under all inertia, so irreducibility makes
wild inertia trivial. The remaining image is cyclic. An
irreducible cyclic subgroup of $\GL_2(\mathbf F_p)$ lies in
$\mathbf F_{p^2}^{\times}$: the minimal polynomial of a
generator is irreducible of degree two. Consequently the
$p$-division field is tame of degree dividing $N$.

By Lemma~\ref{lem:B-RD49-inertia}, a faithful good field for
the actual lattice is also tame of degree dividing $N$.
Let $P$ be their compositum and put $t=[P:K]$, so $t\mid N$.
Choose a uniformizer $\varpi$ of $P$ with
$\varpi^t=\pi_K$, and choose the generator $\sigma$ induced
by $\tau$, so that
\[
                  \sigma(\varpi)=\zeta^{N/t}\varpi.
\]
All $p$-torsion points extend to the good finite flat model,
and lie in its one-dimensional formal group because its
special fibre is connected. Tate full faithfulness extends
the generic descent maps and their inverses to this model
\cite[Theorem~4 and Corollaries~1--2, p.~181]
{Tate1967PDivisible}. Thus the Galois action preserves the
valuation of every nonzero formal point: its coordinate
change has a unit linear term and higher terms have larger
valuation.

The points whose coordinate has valuation at least any
fixed bound form an inertia-stable $\mathbf F_p$-subspace.
If the nonzero points had different valuations, one of
these would be a proper nonzero subspace. They therefore
all have a common positive integer valuation $A$.
The Weierstrass polynomial of $[p](T)$ is monic of degree
$p^2$, has precisely these $p^2$ torsion coordinates as
roots, and differs from $[p](T)$ by an integral unit.
Comparing its linear coefficient with $[p]'(0)=p$ gives
\[
                         NA=e_P=(p-1)t,
 \qquad A=\frac{t}{p+1}.                              \tag{B-RD53.2}
\]
In particular $p+1$ divides $t$.
This is the same one-level product-of-torsion-coordinates
calculation used in Theorem~\ref{thm:B-RD12}; it does not
assume a choice of a CM lift.

Taking the leading coefficient of $T(Q)/\varpi^A$ gives
an injective $\mathbf F_p$-linear map $M[p]\to k$.
The formal group addition has ordinary addition as its
leading term. Under $\sigma$, the coordinate's field action
contributes $\zeta^{(N/t)A}=\zeta^{p-1}$ and the descent
map's linear term contributes $\zeta^c$. Therefore this
two-dimensional $\mathbf F_p$-space is acted on by
multiplication by $\alpha=\zeta^{c+p-1}$.
Irreducibility says $\alpha\notin\mathbf F_p$; its two
eigenvalues are $\alpha,\alpha^p$, proving (B-RD53.1).
The last special-character statement is the two embeddings
of the unramified quadratic subfield of the height-two
division algebra, with the covariant convention of
Proposition~\ref{prop:B-RD4-dictionary}.
\end{proof}

\begin{proposition}[Maximal joint order and adjacent character types]
\label{prop:B-RD53-cycle}
\status{Proved; full independent proof and exact-code audits passed.}
In the remaining $p^2$ comparison of Theorem~\ref{thm:B-RD49},
assume the common residual inertia representation is irreducible
over $\mathbf F_p$. Make the common finite height-one twist
that identifies the common finite determinant with the
level-one tame fundamental character of $K$.
If the two requested special characters differ, their joint
faithful good degree is exactly $N=p^2-1$.

The normalized finite types are the Frobenius orbits of
\[
 c_j=1+(p-1)j\quad\text{in }\mathbf Z/N\mathbf Z,
            \qquad j\in\mathbf Z/(p+1)\mathbf Z.
\]
Label each orbit by its unique even representative $j=2k$,
where $k\in\mathbf Z/n\mathbf Z$ and $n=(p+1)/2$.
The two types in any different-character comparison must
be adjacent vertices in the cycle
\[
                    k\longleftrightarrow k+1
                 \quad\text{on }\mathbf Z/n\mathbf Z. \tag{B-RD53.3}
\]
This is a necessary type condition. The statement does not
construct a pair of groups realizing any edge at level $p^2$.
\end{proposition}
\begin{proof}
The common determinant is nonsquare by
Proposition~\ref{prop:B-RD50-square}. Its quotient by the
level-one tame fundamental character is a square finite
$\mu_{p-1}$-valued character, so the indicated common twist
exists. This character is defined from a uniformizer of $K$;
it is not the mod-$p$ cyclotomic character, which is trivial
because $\mu_p\subset K$. A scalar twist preserves residual
irreducibility.

Let $c_1,c_2\in\mathbf Z/N\mathbf Z$ be the actual tangent
exponents of the two good models. The normalization gives
$c_i\equiv1\pmod{p-1}$. Their isomorphic residual modules
and Lemma~\ref{lem:B-RD53-leading} give either $c_2=c_1$ or
\[
                     c_2=pc_1+(p-1)^2\pmod N.          \tag{B-RD53.4}
\]
The first case has equal special characters and is excluded.
Put $d=\gcd(N,c_1,c_2)$, so the joint faithful order is $N/d$.
Equation (B-RD53.4) implies $d\mid(p-1)^2$, whereas
$c_1\equiv1\pmod{p-1}$ gives $\gcd(d,p-1)=1$.
Thus $d=1$ and the joint order is $N$.

On the $j$-indices, ordinary Frobenius is
$j\mapsto1-j$ modulo $p+1$. Because $p+1$ is even, every
orbit has exactly one even index. The affine involution
(B-RD53.4) is $j\mapsto-1-j$. Applied to the even index
$2k$, it gives an odd index whose Frobenius partner is
$2k+2$. Applied to the other representative of the same
orbit, it gives $2k-2$. Hence the possible distinct-type
relations are exactly the adjacent edges in (B-RD53.3).
No orbit equals its neighbor, since $n\ge3$.
The equality of residual eigenvalue pairs for an abstract
edge follows from the same formula; it does not provide a
Barsotti--Tate realization or a $p^2$-torsion congruence.
\end{proof}

\begin{corollary}[Excluded quadratic pairs and the first new restriction]
\label{cor:B-RD53-quadratic}
\status{Proved; full independent proof and exact-code audits passed.}
For $p\equiv3\pmod4$, $p\ge7$, two different opposite
quadratic special types in the remaining $p^2$ test must
have reducible residual inertia over $\mathbf F_p$.
For $p\equiv1\pmod4$, exactly one unordered quadratic pair
among the normalized abstract finite types is an edge in
(B-RD53.3). Also every remaining normalized joint degree
strictly smaller than $p^2-1$ requires reducible residual
inertia.

For example, at $p=7$ the irreducible-residual possibilities
are the adjacent pairs in the following four-cycle:
\[
 \{1,7\},\quad\{13,43\},\quad\{25,31\},\quad\{19,37\}
                 \quad\text{modulo }48,                \tag{B-RD53.5}
\]
including the edge joining the last pair back to the first.
Its two opposite pairs are excluded in this residual case.
At $p=5$ the cycle is a triangle, so this corollary does
not remove either of the two patterns in
Corollary~\ref{cor:B-RD51-p5}.
\end{corollary}
\begin{proof}
A quadratic twist adds $N/2$ to $c$, hence adds
$n=(p+1)/2$ to $j$. If $n$ is even it acts on cycle
indices by $k\mapsto k+n/2$, which is not an edge for
$n\ge4$. If $n$ is odd it acts by the reflection
\[
                    k\longmapsto(1-n)/2-k\pmod n.
\]
Exactly one unordered exchanged pair is adjacent: solving
$2k=(1-n)/2\mp1$ gives the two ends of the same edge,
because two is invertible modulo odd $n$.
The nonmaximal-order assertion is the preceding proposition.
The displayed four-cycle follows by substituting
$p=7$ and the even indices $0,2,4,6$.
\end{proof}

\begin{remark}[Remaining scope]
\label{rem:B-RD53-scope}
\status{Proved scope and limitation; no least-index or pair-existence assertion.}
The new restriction separates irreducible and reducible
residual inertia over $\mathbf F_p$. It does not assert that
the residual representation of every potentially supersingular
formal group is irreducible. In particular, all remaining
non-CM deformations with reducible residual inertia remain
in the problem. The general interval in
Theorem~\ref{thm:B-RD49} is unchanged.
The retained exact program
\path{code/reduction_descent/verify_tame_residual_cycle.py}
checks $9{,}495$ distinct finite-type pairs and $740$
rational valuation cases for every prime from $5$ through
$101$, including the exact four-cycle at $p=7$, the
triangle at $p=5$, and the quadratic-twist parity assertions.
Its author run passed in $0.04175966899492778$ seconds.
The full independent root and companion audits passed, with
author-output-preserving reruns in $0.04199664699990535$ and
$0.047423224001249764$ seconds, respectively.
The finite calculation supplements the proofs; it does not
instantiate a formal group or assert a torsion-congruent pair.
\end{remark}
